\documentclass[11pt,a4paper]{article}
\usepackage[T1]{fontenc}
\usepackage{lmodern}
\usepackage[margin=26mm]{geometry}
\usepackage{amsmath,amssymb,amsthm,mathtools}
\usepackage{booktabs,array,microtype}
\usepackage{xurl}
\usepackage[colorlinks=true,linkcolor=blue!50!black,urlcolor=blue!50!black,citecolor=blue!50!black]{hyperref}
\usepackage{xcolor}
\hypersetup{pdftitle={An Affine Residual Enclosure for a Projected Resolvent Functional},pdfauthor={Jongmin Choi},pdfsubject={Operator-order residual enclosure and a finite-interval source application}}
\newcommand{\R}{\mathbb R}
\newcommand{\HH}{\mathcal H}
\newcommand{\Id}{\operatorname{Id}}
\newcommand{\Ran}{\operatorname{Ran}}
\newcommand{\Dom}{\operatorname{D}}
\newcommand{\ip}[2]{\langle #1,#2\rangle}
\newcommand{\norm}[1]{\lVert #1\rVert}
\newcommand{\dd}{\,\mathrm dx}
\newcommand{\peta}{p_{\eta}}
\newcommand{\dbar}{\overline d}
\newcommand{\Pp}{\mathcal P_B}
\newtheorem{theorem}{Theorem}[section]
\newtheorem{proposition}[theorem]{Proposition}
\newtheorem{lemma}[theorem]{Lemma}
\newtheorem{corollary}[theorem]{Corollary}
\theoremstyle{remark}
\newtheorem{remark}[theorem]{Remark}
\numberwithin{equation}{section}
\setlength{\emergencystretch}{2em}
\title{An Affine Residual Enclosure\\for a Projected Resolvent Functional}
\author{Jongmin Choi\\[2pt]
\small Independent Researcher, Seoul, Korea\\
\small \href{mailto:24ping@naver.com}{24ping@naver.com}\\
\small ORCID: \href{https://orcid.org/0009-0008-7448-514X}{0009-0008-7448-514X}}
\date{}
\begin{document}
\maketitle
\begin{abstract}
For a positive self-adjoint operator, we record an affine residual enclosure
for a scalar projected resolvent functional. If $A\ge\lambda\Id$, a trial
vector $Ex_0$ has residual $r_0=g-AEx_0$, and $l=Em$, the center contains the
essential correction $\ip{l}{r_0}/(2\lambda)$ and the radius is
$\norm{m}\norm{r_0}/(2\lambda)$. This follows directly from the standard
operator-order estimate $\norm{2\lambda A^{-1}-\Id}\le1$; no claim of
bibliographical priority is made. Exact Galerkin orthogonality removes the
correction, whereas an inexact trial generally does not. We also prove that
the enclosure at $x_0=0$ always contains zero. A concrete application is
specified for an auxiliary operator on a finite interval, with the complete
exterior potential and the actual prime shifts $2,3,4$. Its canonical
weighted graph reserve is reconstructed and distinguished from an
unweighted row sum. An analytic lower bound is proved from explicit
identities and elementary estimates. The required scalar moments and full
residual norms have not been numerically certified for the fixed source.
No sign certification or conclusion about the Riemann hypothesis is obtained.
\end{abstract}

\section{Purpose and scope}
Residual bounds for positive self-adjoint operators are standard consequences
of operator order and spectral calculus; see, for example, the framework in
\cite[Chapter 2]{davies}. Here the purpose is to specify a scalar target,
retain the correction required by an inexact trial, and make the source
definitions of a concrete finite-interval application explicit.

The abstract result is independent of any arithmetic or prime-shift model.
The concrete application uses an auxiliary operator $A_b$ on
$L^2_{\mathrm{even}}((-3/4,3/4);\R)$. It is not the original operator on the
whole real line, and no identification with an original minimum Gamma
realization is asserted. The manuscript contains analytic proofs and a
definition audit. It does not present a newly evaluated response, Gram
matrix, quadrature, or signed numerical enclosure.

\section{The affine residual enclosure}\label{sec:abstract}
Let $\HH$ be a real Hilbert space. Let $A$ be a densely defined self-adjoint
operator with $A\ge\lambda\Id$ for a fixed $\lambda>0$. Thus $A^{-1}$ is a
bounded positive self-adjoint operator on $\HH$. Let
$E:\R^d\to\HH$ be an isometry with $\Ran E\subset\Dom(A)$. Fix
$F\in\HH$, $a\in\R$, and $\alpha,m\in\R^d$, with $E^*F=0$, and define
\begin{equation}\label{eq:abstract-data}
 g=aF-E\alpha,\qquad G=E^*AE,\qquad
 l=Em,\qquad M=\ip{l}{A^{-1}g}.
\end{equation}
The adjoint $E^*$ and all norms refer to the Hilbert structure in question.
For any fixed $x_0\in\R^d$, set
\begin{equation}\label{eq:residual}
 r_0=g-AEx_0,\qquad
 \delta=m^T(-\alpha-Gx_0).
\end{equation}
Because $E^*g=-\alpha$, one has $E^*r_0=-\alpha-Gx_0$ and
$\delta=\ip{Em}{r_0}$.

\begin{theorem}[Affine residual enclosure]\label{thm:affine}
Under these assumptions,
\begin{equation}\label{eq:main}
 \boxed{\left|M-m^Tx_0-\frac{\delta}{2\lambda}\right|
 \le\frac{\norm{m}\norm{r_0}}{2\lambda}.}
\end{equation}
In particular the residual norm in this statement is the full norm in
$\HH$, not the norm of its projection $E^*r_0$.
\end{theorem}
\begin{proof}
The spectral theorem gives
$0\le A^{-1}\le\lambda^{-1}\Id$, and hence
\begin{equation}\label{eq:contraction}
 \norm{2\lambda A^{-1}-\Id}\le1,
 \qquad
 \norm{A^{-1}-(2\lambda)^{-1}\Id}\le(2\lambda)^{-1}.
\end{equation}
Since $Ex_0\in\Dom(A)$,
$A^{-1}g=Ex_0+A^{-1}r_0$. The isometry of $E$ therefore yields
\[
 M-m^Tx_0-\frac{\delta}{2\lambda}
 =\ip{Em}{\bigl(A^{-1}-(2\lambda)^{-1}\Id\bigr)r_0}.
\]
Apply Cauchy--Schwarz and \eqref{eq:contraction}, using
$\norm{Em}=\norm{m}$.
\end{proof}

\begin{corollary}[Exact Galerkin trial]\label{cor:galerkin}
The matrix $G$ is symmetric and $G\ge\lambda I_d$. The unique solution
$x_*=-G^{-1}\alpha$ satisfies $E^*(g-AEx_*)=0$. Thus
\begin{equation}\label{eq:galerkin}
 \left|M-m^Tx_*\right|
 \le\frac{\norm{m}\norm{g-AEx_*}}{2\lambda}.
\end{equation}
For a general inexact $x_0$, the correction in \eqref{eq:main} must be retained.
\end{corollary}
\begin{proof}
For $y\in\R^d$,
$y^TGy=\ip{Ey}{AEy}\ge\lambda\norm{y}^2$. Symmetry follows from
self-adjointness and $\Ran E\subset\Dom(A)$. The remaining assertions
follow from $E^*r_0=-\alpha-Gx_0$.
\end{proof}

\begin{remark}[Comparison with the usual residual bound]\label{rem:comparison}
The ordinary estimate is
$|M-m^Tx_0|\le\norm{m}\norm{r_0}/\lambda$.
For the same full residual and the same $\lambda$, \eqref{eq:main} has half
that radius, but generally a different center. Its exact interval is
contained in the ordinary interval because
$|\delta|\le\norm{m}\norm{r_0}$. If $\delta$, $m$, or $r_0$ are themselves
enclosed imperfectly, their additional errors must also be included. No
factor-of-two improvement in an actual stored numerical certificate follows
from this comparison alone. The constant in \eqref{eq:contraction} is the
sharp uniform constant, already attained by $A=\lambda\Id$.
\end{remark}

\begin{proposition}[Zero-start limitation]\label{prop:zero-start}
For $x_0=0$, the norm-based interval in \eqref{eq:main} always contains zero.
Consequently this interval alone cannot certify a strict sign of $M$ from
the zero trial.
\end{proposition}
\begin{proof}
Here $r_0=g$, $m^Tx_0=0$, and $\delta=\ip{Em}{g}$. Cauchy--Schwarz gives
$|\delta|\le\norm{m}\norm{g}$. The absolute value of the center
$\delta/(2\lambda)$ is therefore no larger than the radius
$\norm{m}\norm{g}/(2\lambda)$.
\end{proof}
The proposition does not assert $M=0$. To obtain a sign by this mechanism,
one needs a nonzero useful trial or additional information beyond this
general norm interval. It supplies no reason to start a zero-trial numerical
experiment merely to certify a sign.

\section{The finite-interval source operator}\label{sec:source}
Fix
\begin{equation}\label{eq:constants}
 B=\frac34,\quad \Omega=(-B,B),\quad
 \theta=\frac1{65536},\quad\omega=1-\theta,\quad
 a=\frac{261572074620463589}{296898235056783360}>0.
\end{equation}
Let $t=\log2$, $u=\log3$, and
\[
 c_2=\frac{t}{\sqrt2},\qquad c_3=\frac{u}{\sqrt3},\qquad c_4=\frac t2.
\]
Every function on $\Omega$ is extended by zero outside it. If
$T_s f(x)=\widetilde f(x-s)$ for $x\in\Omega$, define
\begin{equation}\label{eq:prime}
 \Pp=c_2(T_t+T_{-t})+c_3(T_u+T_{-u})
                  +c_4(T_{2t}+T_{-2t}).
\end{equation}
This is a bounded self-adjoint operator, preserves evenness, and retains
both directions of the shifts. The coefficient of the shift $\log4=2t$ is
$\log2/2$, with no additional duplicate prime-power term.

For $s>0$ put
\begin{align}
 \rho(s)&=\frac{e^{-s/2}}{1-e^{-2s}},\qquad
 R(s)=\int_s^\infty\rho(v)\,\mathrm dv
       =\operatorname{artanh}(e^{-s/2})+\arctan(e^{-s/2}),\label{eq:tail}\\
 V(x)&=R(B+x)+R(B-x),\qquad V_0=2R(B),\label{eq:potential}\\
 c_\Gamma&=\log\pi+\gamma_E+\frac\pi2+3\log2,\qquad
 \kappa_g=-c_\Gamma+\frac{17}{5}\omega+7\theta.\label{eq:kappa}
\end{align}
Here $\gamma_E$ is Euler's constant. The tail in \eqref{eq:tail} is complete.
The two terms in \eqref{eq:potential} are the two exterior costs associated
with zero extension; neither is truncated or multiplied by two again.

On $\HH_b=L^2_{\mathrm{even}}(\Omega;\R)$ define the auxiliary operator
\begin{equation}\label{eq:Ab}
 A_b=(a+\kappa_g)\Id+\omega M_V-\Pp,\qquad
 \Dom(A_b)=\{f\in\HH_b:Vf\in L^2(\Omega)\},
\end{equation}
where $M_V$ is multiplication by $V$. Its form domain is
\begin{equation}\label{eq:form-domain}
 \mathcal V_b=\left\{f\in\HH_b:\int_\Omega V|f|^2\dd<\infty\right\}.
\end{equation}
The real multiplication operator, plus a bounded self-adjoint perturbation,
is self-adjoint on the displayed operator domain. Its semibounded closed
form has domain \eqref{eq:form-domain}. The positive lower bound will be
proved in Section~\ref{sec:lower}. In \eqref{eq:Ab}, $\omega$ multiplies
the potential; the off-diagonal prime coefficients are $-c_2,-c_3,-c_4$,
not $-\omega c_j$ and not $-2c_j$. The constant $\kappa_g$ also depends on
$\omega$ through \eqref{eq:kappa}.

\subsection{Physical basis and the fixed polynomial data}
Use the physically orthonormal even Legendre functions
\begin{equation}\label{eq:basis}
 e_{2j}(x)=\sqrt{\frac{4j+1}{2B}}P_{2j}(x/B),\qquad0\le j<16,
\end{equation}
and let $Ey=\sum_{j=0}^{15}y_j e_{2j}$ and $P=EE^*$. The normalization
follows from Legendre orthogonality; see \cite[Section 18.3]{dlmf}.
The fixed high polynomial is specified by rational coefficients $b_m$ as
\begin{equation}\label{eq:F}
 F(x)=\frac1{\sqrt{2B}}\sum_{m=32,34,\ldots,128}b_m P_m(x/B).
\end{equation}
There are 49 coefficients. The exact fixed instance and the rational vector
$\alpha\in\R^{16}$ are preserved as source strings in
\nolinkurl{source_definitions.json}; no decimal endpoint transcription is
used to define them. The results here hold for any such fixed coefficients.
The archived instance is identified by the source hashes in
Section~\ref{sec:availability}, and is not asserted to be an exact full
inverse merely because it arose as a finite inverse prefix.

The degree separation gives $E^*F=0$. The endpoint growth
$V(x)=O(1+|\log(B-|x|)|)$ follows from \eqref{eq:tail}. Since
$\int_0^\varepsilon|\log s|^2\,\mathrm ds<\infty$, every $e_{2j}$ and
$F$ belongs to $\Dom(A_b)$. This establishes the operator membership
needed for $g-A_bEx_0$, which is stronger than mere form membership.
It does not prove membership in any original whole-line minimum operator
domain.

\section{The canonical weighted graph reserve}\label{sec:reserve}
\subsection{Complete component geometry}
Put $r_g=B-t$, $v_g=u-t$, and $w_g=2t-u$. Up to null endpoints, the graph
of the shifts in \eqref{eq:prime} is described by three component charts:
\begin{align*}
 J_7&=(-r_g,r_g),&
 \xi_7(x)&=(x,x+t,x-t,x+v_g,x-v_g,x+w_g,x-w_g),\\
 J_4&=(u+B-2t,2t-B),&
 \xi_4(y)&=(y-t,y,y-u,y-u+t),\\
 J_2&=(B+t-u,u-B),&
 \xi_2(z)&=(z-t,z).
\end{align*}
All coordinate Jacobians are one. With zero-based vertex labels, the edges
are
\begin{center}
\small
\begin{tabular}{clll}
\toprule
Chart&$c_2$ edges&$c_3$ edges&$c_4$ edges\\
\midrule
$7$&$(0,1),(0,2),(4,5),(3,6)$&$(1,4),(2,3)$&$(1,2)$\\
$4$&$(0,1),(2,3)$&$(1,2)$&none\\
$2$&$(0,1)$&none&none\\
\bottomrule
\end{tabular}
\end{center}

For completeness, the positive physical half is partitioned by the seven
bands in Table~\ref{tab:reserve}. Their vertex images are disjoint except
at endpoints. The negative half is obtained by reflection: the chart
parameters map as $x\mapsto-x$, $y\mapsto u-y$, and $z\mapsto t-z$, with
the corresponding vertex permutations. Checking the translated coordinates
against $\pm B$ gives exactly the edge table above. In particular both
central $n=2$ neighbors and the $n=4$ edge occur.

The inequalities needed for this partition follow from the positive series
\begin{equation}\label{eq:log-series}
 2\sum_{j=0}^{N-1}\frac{q^{2j+1}}{2j+1}
 <\log\frac{1+q}{1-q}
 <2\sum_{j=0}^{N-1}\frac{q^{2j+1}}{2j+1}
       +\frac{2q^{2N+1}}{(2N+1)(1-q^2)},\quad0<q<1.
\end{equation}
Using $(q,N)=(1/3,3),(1/2,4)$ gives
$693/1000<t<694/1000$ and $1098/1000<u<1099/1000$.
These bounds imply $0<r_g<w_g<v_g<t<B$,
$2u-t>2B$, $4t-u>2B$, $t+u>2B$, and $3t-u>B$.
Also $\log5>3/2$, from the first three positive terms in
\eqref{eq:log-series} with $q=2/3$. Thus shifts with $n\ge5$ have no
positive-measure overlap on this support.

Translation of representatives therefore gives the unitary identification
\[
 L^2(\Omega)\simeq L^2(J_7;\R^7)\oplus
 L^2(J_4;\R^4)\oplus L^2(J_2;\R^2).
\]
This is an assertion about almost-everywhere representatives, not point
evaluation of an abstract $L^2$ class. It restricts to the reflection-even
reducing subspace. Each undirected edge contributes
$2c_{ij}f_i f_j$, representing both arithmetic directions once.

\subsection{Weighted rows and the reserve function}
For positive weights $w_i$, the exact edge identity is
\begin{equation}\label{eq:edge-square}
 c_{ij}\left(\frac{w_j}{w_i}|z_i|^2+
              \frac{w_i}{w_j}|z_j|^2-2z_i z_j\right)
 =c_{ij}w_iw_j\left|\frac{z_i}{w_i}-\frac{z_j}{w_j}\right|^2.
\end{equation}
Consequently the diagonal associated with these squares is the
\emph{weighted} actual row
$d_{\mathrm{act},i}=\sum_{j\sim i}c_{ij}w_j/w_i$.

The elementary bounds $c_2<1/2$, $c_3<13/20$, $c_4<7/20$ follow from
the logarithm bounds above and $1/\sqrt2<71/100$, $1/\sqrt3<7/12$.
They are used only for inequalities, not to redefine \eqref{eq:prime}.
Choose reserve weights $(10/9,1,9/10,1/2)$ for the seven chart's
center, $t$, $v$, and $w$ vertex types. The four weighted row bounds are
\[
 \frac9{10},\qquad
 \frac59+\frac{117}{200}+\frac7{20}=\frac{2683}{1800},\qquad
 \frac{13}{18}+\frac5{18}=1,\qquad\frac9{10}.
\]
For the path choose weights $3/5$ at leaves and $1$ at interior vertices;
the bounds are $5/6$ and $19/20$. For the pair use equal weights and the
bound $1/2$. This defines the bounded even function $\dbar$ in the table.

\begin{table}[htbp]
\centering\small
\begin{tabular}{lll}
\toprule
Positive physical band&Vertex type&$\dbar$\\
\midrule
$(0,r_g)$&seven center&$9/10$\\
$(r_g,3t-u-B)$&path leaf&$5/6$\\
$(3t-u-B,B+t-u)$&seven $w$&$9/10$\\
$(B+t-u,u-B)$&pair&$1/2$\\
$(u-B,u+B-2t)$&seven $v$&$1$\\
$(u+B-2t,2t-B)$&path interior&$19/20$\\
$(2t-B,B)$&seven $t$&$2683/1800$\\
\bottomrule
\end{tabular}
\caption{The canonical reserve on $\Omega$. Use even reflection on the
negative half; values at the finitely many boundaries are immaterial.}
\label{tab:reserve}
\end{table}

Summing \eqref{eq:edge-square} gives, on finite weighted energy vectors,
\begin{equation}\label{eq:Sgraph}
 S_{\mathrm{graph}}[f]
 :=\int_\Omega\dbar|f|^2\dd-\ip{f}{\Pp f}\ge0.
\end{equation}
Its exact decomposition retains the weighted edge squares and the
nonnegative diagonal slack $\dbar-d_{\mathrm{act}}$.

\begin{proposition}[The reserve is not an unweighted row sum]\label{prop:correction}
The definition $\dbar=\Pp1_\Omega$ is incompatible with
Table~\ref{tab:reserve}. It changes the multiplier
$\peta=\omega(W_B+9/10-\dbar)$ and any moment derived from it.
In fact that replacement makes this multiplier negative near the center.
\end{proposition}
\begin{proof}
On the central band only the two $n=2$ shifts meet the constant supported
function, so $\Pp1_\Omega=2c_2=\sqrt2\log2$, whereas the canonical
reserve is $9/10$. Since $\log2>2/3$ and $\sqrt2>7/5$,
$\sqrt2\log2>14/15>9/10$. At zero, $W_B(0)=0$, and continuity on the
central band proves the asserted negativity on a positive-measure
neighborhood.
\end{proof}
The distinction also applies to $d_{\mathrm{act}}$, which is weighted,
not $\Pp1_\Omega$. Replacing $\dbar$ consistently in an algebraic
decomposition does not alter \eqref{eq:Ab}, but it alters $\peta$, $m$, and
the scalar target. A bounded even change of $\dbar$ leaves the relevant
weighted domain equivalence intact; it does not preserve $\peta\ge0$ or
the numerical functional.

\section{An explicit analytic lower bound}\label{sec:lower}
We reconstruct the estimates required for a lower bound. No stored audit
status or numerical response output is used as a proof premise.

\begin{lemma}[Complete logarithmic reserve]\label{lem:log-reserve}
For $x\in\Omega$,
\[
 V(x)\ge V_0+W_B(x),\qquad
 W_B(x)=-\frac12\log(1-x^2/B^2).
\]
Moreover $\eta=W_B+9/10-\dbar\ge0$ almost everywhere.
\end{lemma}
\begin{proof}
Write $k(s)=\rho(s)-1/(2s)$ and $z=s/2$. Direct algebra gives
$k(s)=[\operatorname{csch}z-1/z+\operatorname{sech}z]/4$.
For $0<z\le3/4$, the power series imply
\[
 \frac{\sinh z}{z}\le\frac1{1-z^2/6},\qquad
 \cosh z\ge1+z^2/2.
\]
The first inequality follows from $(2j+1)!\ge6^j$. For $b=z^2\le9/16$,
\[
 (1+b/2)(1-b/6)^2-1=\frac{b(12-10b+b^2)}{72}>0.
\]
Hence $z^2\cosh z>\sinh^2z$, so
$\operatorname{csch}z-1/z$ is decreasing; so is $\operatorname{sech}z$.
Thus $k$ decreases for $0<s\le3/2$. Since $R'=-\rho$, for $0<x<B$,
\[
 V'(x)=\rho(B-x)-\rho(B+x)
 \ge\frac1{2(B-x)}-\frac1{2(B+x)}
 =\frac{x}{B^2-x^2}=W_B'(x).
\]
Integration from zero and evenness prove the first assertion.

Only the large reserve rows require additional estimates. Since
$t>69/100$ and $u>109/100$, \eqref{eq:log-series} gives
\begin{align*}
 W_B(2t-B)&>\tfrac12\log(625/184)
 >\frac{441}{809}+\frac13\left(\frac{441}{809}\right)^3
 >\frac{1063}{1800},\\
 W_B(u-B)&>\tfrac12(34/75)^2=\frac{578}{5625}>\frac1{10}.
\end{align*}
The last finite comparisons are rational inequalities after clearing
positive denominators. Seven $t$ vertices have $|x|\ge2t-B$, and seven
$v$ vertices and path interior vertices have $|x|\ge u-B$.
The values of $\eta$ are consequently
\[
\begin{array}{c|c}
\text{vertex type}&\eta\\\hline
\text{seven center or }w&W_B\\
\text{seven }t&W_B-1063/1800\\
\text{seven }v&W_B-1/10\\
\text{path leaf}&W_B+1/15\\
\text{path interior}&W_B-1/20\\
\text{pair}&W_B+2/5.
\end{array}
\]
Each is nonnegative on its physical band.
\end{proof}

\begin{lemma}[Prime norm bound]\label{lem:prime-norm}
The actual operator satisfies $\norm{\Pp}\le6/5$.
\end{lemma}
\begin{proof}
These weights are different from the reserve weights. On the seven chart
choose $(5/6,1,2/3,5/18)$ for the center, $t$, $v$, and $w$ types.
The coefficient bounds already proved give the following neighbor sums:
\[
\begin{array}{c|c|c}
\text{type}&\sum_{j\sim i}c_{ij}w_j\text{ upper bound}&(6/5)w_i\\\hline
\text{center}&1&1\\
t&6/5&6/5\\
v&71/90&4/5\\
w&1/3&1/3.
\end{array}
\]
Apply $2c_{ij}|z_i z_j|\le
c_{ij}(w_j/w_i)|z_i|^2+c_{ij}(w_i/w_j)|z_j|^2$ to each edge.
For the path, equal weights give a maximal row below $23/20<6/5$;
for the pair it is below $1/2$. The direct integral yields
$|\ip{f}{\Pp f}|\le(6/5)\norm{f}^2$ and hence the norm bound.
\end{proof}

\begin{lemma}[A rational exterior surplus]\label{lem:delta}
The exact constants satisfy
\begin{equation}\label{eq:delta-bound}
 \delta_0:=V_0-c_\Gamma+\frac52>\frac{29}{14000},\qquad
 j_7:=\frac{29}{5}-c_\Gamma>0.
\end{equation}
\end{lemma}
\begin{proof}
With $z=e^{-3/8}$, the two inverse trigonometric series combine to give
$V_0=4\sum_{j\ge0}z^{4j+1}/(4j+1)$.
The alternating exponential polynomial through degree five gives
\[
 z>\sum_{j=0}^{5}\frac{(-3/8)^j}{j!}
   =\frac{900839}{1310720}>\frac{687}{1000}=:b.
\]
The omitted terms in the combined series are positive, and the finite
rational inequality
\[
 4\left(b+\frac{b^5}{5}+\frac{b^9}{9}\right)>\frac{5771}{2000}
\]
gives $V_0>5771/2000$.

For the upper constant, the positive identity
$22/7-\pi=\int_0^1 x^4(1-x)^4/(1+x^2)\,\mathrm dx$ gives $\pi<22/7$.
Since $\sum_{j=0}^{5}(23/20)^j/j!>22/7$, also $\log\pi<23/20$.
Equation~\eqref{eq:log-series} gives
\[
 \frac{842}{1215}<\log2<\frac{23581}{34020}<\frac{347}{500}.
\]
Writing $H_n=\sum_{k=1}^n1/k$, a telescoping integral yields
\[
 H_n-\log n-\gamma_E
 =\sum_{k=n}^\infty\int_0^1
   \frac{1-y}{(k+y)(k+1)}\,\mathrm dy
 >\frac12\sum_{k=n}^\infty(k+1)^{-2}>\frac1{2(n+1)}.
\]
As $H_{16}=2436559/720720$, it follows that
\[
 \gamma_E<\frac{2436559}{720720}
              -4\frac{842}{1215}-\frac1{34}<\frac{29}{50}.
\]
These finite rational inequalities are verified by clearing denominators;
their formulas, rather than a prior computation's status, are the evidence.
Consequently
\[
 c_\Gamma<\frac{23}{20}+\frac{29}{50}+\frac{11}{7}
                         +3\frac{347}{500}=\frac{9421}{1750},
\]
and
\[
 \delta_0>\frac{5771}{2000}-\frac{9421}{1750}+\frac52
           =\frac{29}{14000}.
\]
The same upper bound for $c_\Gamma$ gives $j_7>0$.
\end{proof}

\begin{proposition}[Lower bound for the same auxiliary operator]\label{prop:lower}
For the operator \eqref{eq:Ab},
\begin{equation}\label{eq:lambda-b}
 A_b\ge\lambda_b\Id,\qquad
 \lambda_b=a+\frac{29\omega}{14000}>0.
\end{equation}
The bound concerns this auxiliary finite-interval operator only.
\end{proposition}
\begin{proof}
Put $\Delta_\Gamma=V-V_0-W_B\ge0$ and
$J_7[f]=(7-c_\Gamma)\norm{f}^2-\ip{f}{\Pp f}$.
Lemmas~\ref{lem:prime-norm} and \ref{lem:delta} imply
$J_7[f]\ge j_7\norm{f}^2$.
The closed form $\mathfrak a_b$ of \eqref{eq:Ab} has the exact decomposition
\begin{align}\label{eq:decomposition}
 \mathfrak a_b[f]={}&a\norm{f}^2+\int_\Omega\peta|f|^2\dd
       +\omega S_{\mathrm{graph}}[f]\\
 &+\omega\int_\Omega(\Delta_\Gamma+\delta_0)|f|^2\dd
       +\theta J_7[f],\nonumber\\
 \peta={}&\omega(W_B+9/10-\dbar).\nonumber
\end{align}
Indeed $\dbar$ and $W_B$ cancel, and $\omega+\theta=1$ leaves the
arithmetic term $-\ip{f}{\Pp f}$ exactly once. The remaining constant is
$-c_\Gamma+(17/5)\omega+7\theta=\kappa_g$.
Each nonconstant term in \eqref{eq:decomposition} is nonnegative by the
preceding lemmas. Thus
\[
 \mathfrak a_b[f]\ge
 (a+\omega\delta_0+\theta j_7)\norm{f}^2
 \ge\lambda_b\norm{f}^2.
\]
The bound extends over the closed form domain and gives the associated
operator inequality. The contribution of $\delta_0$ occurs once.
\end{proof}

The residual reserve relative to the multiplier
$q_0=a+\peta+\theta j_7$ is at least $29\omega/14000$.
This is not a computed smallest eigenvalue. No separate full-high diagonal
reserve, internal-kernel remainder, or whole-pole term has been added to
$A_b$ to obtain \eqref{eq:lambda-b}. The complete pole in a different trial
objective remains a separate object.

\section{The projected functional and the missing numerical suppliers}\label{sec:application}
For the canonical source define
\begin{equation}\label{eq:source-M}
 g=aF-E\alpha,\quad c_b=A_b^{-1}g,\quad
 m=E^*(\peta F),\quad
 M_b=\ip{\peta F}{Pc_b}=\ip{Em}{A_b^{-1}g}.
\end{equation}
The function $\peta F$ belongs to $L^2$ by logarithmic endpoint growth.
Thus Theorem~\ref{thm:affine} applies with $\lambda=\lambda_b$ and
$G=E^*A_bE$. It also applies with the weaker valid lower bound $\lambda=a$.
For a fixed $x_0$, define
\[
 C=m^Tx_0,\quad D=m^T(-\alpha-Gx_0),\quad
 N_m\ge\norm{m},\quad N_r\ge\norm{g-A_bEx_0}.
\]
Certified full-domain enclosures of these scalars would give
\begin{equation}\label{eq:scalar-certificate}
 M_b\in C+\frac{D}{2\lambda}
              +\left[-\frac{N_mN_r}{2\lambda},\frac{N_mN_r}{2\lambda}\right],
\end{equation}
with outward treatment of every input error. For a single fixed $x_0$,
one need not store the whole matrix $G$: a certified $Gx_0$ or the direct
scalar $D$ suffices, together with the other suppliers.

\subsection{Quantities that cannot substitute for these inputs}
In the audited stored files, the vector named \texttt{chi} has target
$E^*A_bF$, not $m=E^*(\peta F)$. The scalar named \texttt{H} in the form
packet has target $\mathfrak a_b[F]+2|\ell|^2$ for a separate whole-pole
moment $\ell$; it is not the $16\times16$ matrix $G$.
The stored \texttt{beta} vector is a response projection $E^*c_b$; it is
neither $m$ nor a prescribed certified $x_0$. Fibre response centers of
dimensions $7,4,2$ are not one global low-dimensional trial $Ex_0$.
Other residual or observation Grams have different operators and bases.

In particular
\[
 \norm{E^*r_0}=\norm{-\alpha-Gx_0}\le\norm{r_0}
\]
is a lower bound on the full residual norm, not a replacement upper bound.
The omitted high component $(\Id-P)r_0$ matters. An energy bound
$G\le C_{\mathrm{low}}I$ also does not imply
$\norm{A_bEx_0}\le C_{\mathrm{low}}\norm{x_0}$.

\subsection{Scope of the data conclusion}
Within the files identified in the accompanying source manifest, no
certified $m$, $Gx_0$, full residual norm, or matching scalar suppliers
$C,D,N_m,N_r$ were found. The conclusion is limited to that inspection
scope, not all other workspaces or possible data holdings. The earlier
full independent response verification was incomplete; its stored
projection intervals are not used here as certified numerical suppliers.
The numerical result affected by transcribed beta endpoints is also
excluded. No new source response or Gram calculation was performed to
fill the gap.

The exact polynomial and coefficient data define the target; their byte
identity does not establish the truth of a numerical enclosure. The analytic
lower bound in Proposition~\ref{prop:lower} is proved above from the exact
operator definition. It does not turn an uncertified response output into
a certificate. At present \eqref{eq:scalar-certificate} is a specified
certificate format, not an evaluated sign result.

\section{Data availability}\label{sec:availability}
The source definitions, reproducibility scripts and audit corrections are
publicly available at the immutable source/data revision
\begin{center}
\url{https://github.com/serabing-hash/theta-weighted-certificate/tree/ed9ce2c814e2ef64da4ea0337ada8257bb8b0c8f/research/2026-10-13-direct-m-residual-enclosure-paper/}.
\end{center}
The source/data commit is
\texttt{ed9ce2c814e2ef64da4ea0337ada8257bb8b0c8f}.
Its manifest and files were read back through anonymous public GitHub URLs
and matched the local bytes before the final manuscript was compiled.
This fixed link identifies the source/data deposit; it does not need to
identify the later commit containing this PDF itself.

The final standalone TeX, PDF and a ZIP containing exactly those two files
are deposited separately in the same repository directory. The repository is
\url{https://github.com/serabing-hash/theta-weighted-certificate}.
The source/data revision contains \nolinkurl{source_definitions.json},
the scientific input manifest, the reproduction README, audit corrections,
and the document/static verification scripts. Internal filesystem paths,
Library identifiers and private transfer URLs are not part of this deposit.

Exact $b_m$ and $\alpha$ strings are copied mechanically from the audited
inputs, with scientific source hashes retained. The fixed-polynomial and
alpha input hashes have prefixes \texttt{0d702949feba1632} and
\texttt{7fb2fede946e3899}; the input manifest contains their complete hashes.
The recovered source-definition text has hash prefix
\texttt{e31b29a48d2b4dcd}. Historical source archives are not redistributed
in this directory. The required analytic definitions and proofs are written
in this manuscript, while the exact fixed coefficient strings are supplied
with the public definitions. Hashes establish data identity, not proof by
an earlier audit's PASS status.

The TeX file includes its bibliography internally and requires no external
figure, bibliography, or coefficient input to compile. No earlier paper by
the author is cited. The proof scope is the abstract residual estimate and
the explicitly defined finite-interval auxiliary operator. A relationship
to an original whole-line realization, a numerical sign for $M_b$, and any
Riemann-hypothesis conclusion remain outside the established results.

\begin{thebibliography}{9}
\bibitem{davies}
E. B. Davies,
\emph{Spectral Theory and Differential Operators},
Cambridge Studies in Advanced Mathematics 42,
Cambridge University Press, 1995.
\href{https://doi.org/10.1017/CBO9780511623721}{doi:10.1017/CBO9780511623721}.
The spectral theorem is developed in Chapter 2, pp.~22--44.

\bibitem{dlmf}
National Institute of Standards and Technology,
\emph{Digital Library of Mathematical Functions},
Section 18.3, ``Definitions'' (classical orthogonal polynomials),
\url{https://dlmf.nist.gov/18.3}.
\end{thebibliography}
\end{document}
